\documentclass[11pt,a4paper]{article}

\usepackage[utf8]{inputenc}
\usepackage[T1]{fontenc}
\usepackage[turkish]{babel}
\usepackage{lmodern}
\usepackage{geometry}
\usepackage{microtype}
\usepackage{booktabs}
\usepackage{tabularx}
\usepackage{longtable}
\usepackage{array}
\usepackage{enumitem}
\usepackage{xcolor}
\usepackage{hyperref}
\usepackage{fancyhdr}
\usepackage{titlesec}

\geometry{top=2.2cm,bottom=2.2cm,left=2.2cm,right=2.2cm}
\hypersetup{
  colorlinks=true,
  linkcolor=blue!45!black,
  urlcolor=blue!55!black,
  pdftitle={İHA Telemetrisinde Anomali Tespiti Proje Süreci ve Fizibilite Raporu},
  pdfauthor={Anıl Keskin}
}

\definecolor{navy}{HTML}{17365D}
\definecolor{softblue}{HTML}{EAF2F8}
\definecolor{softgray}{HTML}{F3F4F6}
\definecolor{darkgreen}{HTML}{1B5E20}
\definecolor{darkred}{HTML}{8B1E1E}

\titleformat{\section}{\Large\bfseries\color{navy}}{\thesection.}{0.6em}{}
\titleformat{\subsection}{\large\bfseries\color{navy}}{\thesubsection.}{0.6em}{}
\setlist[itemize]{leftmargin=1.6em,itemsep=3pt,topsep=4pt}
\setlist[enumerate]{leftmargin=1.8em,itemsep=4pt,topsep=4pt}
\renewcommand{\arraystretch}{1.18}

\pagestyle{fancy}
\fancyhf{}
\lhead{\small İHA Telemetrisi Anomali Tespiti}
\rhead{\small Proje Süreci ve Fizibilite}
\cfoot{\thepage}

\newcommand{\code}[1]{\texttt{\detokenize{#1}}}
\newcommand{\resultok}[1]{\textcolor{darkgreen}{\textbf{#1}}}
\newcommand{\resultno}[1]{\textcolor{darkred}{\textbf{#1}}}

\begin{document}

\begin{titlepage}
  \centering
  \vspace*{2.4cm}
  {\Huge\bfseries\color{navy}
  İHA Telemetrisinde\\[5pt]
  Anomali Tespiti\par}
  \vspace{0.7cm}
  {\LARGE Proje Süreci, Teknik Kararlar ve\\
  Fizibilite Değerlendirmesi\par}
  \vspace{1.3cm}
  {\large Yönetici ve teknik değerlendirme görüşmeleri için hazırlanmıştır\par}
  \vfill

  \begin{tabular}{rl}
    \textbf{Hazırlayan:} & Anıl Keskin \\
    \textbf{Tarih:} & 16 Temmuz 2026 \\
    \textbf{Kapsam:} & Veri hazırlama, özellik mühendisliği, ML/DL, fiziksel yöntemler, doğrulama \\
  \end{tabular}

  \vfill
  \fcolorbox{navy}{softblue}{
    \begin{minipage}{0.88\textwidth}
      \textbf{Raporun amacı:} Çalışmayı depo içi deney kodlarıyla değil; hedef,
      kullanılan uçuş telemetrisi, karşılaşılan problemler, uygulanan çözümler,
      yöntem değişikliklerinin gerekçesi ve nihai teknik hüküm üzerinden anlatmaktır.
    \end{minipage}}
  \vspace*{1.4cm}
\end{titlepage}

\tableofcontents
\newpage

\section{Yönetici özeti}

Bu proje tarafıma, İHA ve uçuş telemetrisi üzerinde anomali tespitinin makine
öğrenmesiyle ne ölçüde gerçekleştirilebileceğinin araştırılması amacıyla verildi.
Veri kaynaklarının araştırılması ve indirilmesi, farklı kayıt biçimleri için
okuyucuların geliştirilmesi, veri temizleme, özellik üretimi, modelleme, deney
tasarımı, hata analizi ve teknik raporlama çalışmalarının büyük bölümü tarafımdan
gerçekleştirildi.

Çalışma yalnızca ``bir model eğitip doğruluk değerine bakma'' şeklinde
yürütülmedi. Beş farklı veri ailesi üzerinde gerçek arıza, siber saldırı, simülasyon
ve etiketsiz doğal uçuş davranışı birlikte incelendi. Klasik makine öğrenmesi,
derin öğrenme, istatistiksel değişim tespiti ve uçuş fiziğine dayalı kontrol
yöntemleri karşılaştırıldı. LSTM Autoencoder, Dense Autoencoder, USAD, Isolation
Forest, LightGBM ve bağlamsal LSTM gibi öğrenen yöntemlerin yanında CUSUM,
robust z-skoru, otopilot kalıntıları ve fiziksel tutarlılık kontrolleri kullanıldı.

Süreç boyunca ilk bakışta başarılı görünen bazı sonuçların veri sızıntısı, hatalı
etiket kapsamı, ölçek/genlik etkisi veya aynı uçuşun eğitim ve testte temsil
edilmesi gibi nedenlerle gerçeği olduğundan iyi gösterdiği tespit edildi. Bu
problemler düzeltildikçe sonuçların düşmesi bir başarısızlık değil, çalışmanın
güvenilir hâle gelmesidir.

Nihai bulgu şudur:

\begin{quote}
\resultno{Mevcut veri, etiket ve sensör gözlenebilirliği ile farklı İHA
platformlarına genellenebilen, kabul edilebilir arıza yakalama oranını düşük
yanlış alarm yüküyle birlikte sağlayan genel amaçlı bir anomali tespit ürünü
teknik olarak doğrulanamamıştır.}
\end{quote}

Bu hüküm ``makine öğrenmesi hiçbir şey öğrenemedi'' anlamına gelmemektedir.
Belirli veri setlerinde ve belirli arıza türlerinde güçlü ayrışma elde edilmiştir.
Ancak model eşiği arızaları yakalayacak kadar gevşetildiğinde operatörün
kullanamayacağı kadar çok alarm oluşmuş; yanlış alarmları azaltacak kadar
sıkılaştırıldığında ise gerçek arızaların önemli bölümü kaçırılmıştır. Son yapılan
GNSS bütünlük çalışması da aynı temel sınırı daha kontrollü bir deneyle
doğrulamıştır.

\section{Projenin başlangıç hedefi ve benim sorumluluğum}

\subsection{Başlangıçta hedeflenen}

İlk hedef; İHA telemetrisindeki olağan dışı davranışları, her arıza için ayrı bir
el yazımı kural gerektirmeden tespit edebilecek bir makine öğrenmesi sistemi
geliştirmekti. Beklenen sistemin:

\begin{itemize}
  \item normal uçuş davranışını öğrenmesi,
  \item sensör, aktüatör, navigasyon veya kontrol kaynaklı sapmaları fark etmesi,
  \item mümkünse farklı uçuş ve platformlarda çalışması,
  \item gerçek arızaları yeterli oranda yakalarken düşük yanlış alarm üretmesi
\end{itemize}

amaçlandı.

\subsection{Fiilen yürüttüğüm çalışmalar}

Projenin teknik sahipliği kapsamında aşağıdaki işler uçtan uca ele alındı:

\begin{itemize}
  \item açık ve kurum içi erişilebilir veri kaynaklarının araştırılması,
  \item ROS bag, PX4 ULog, CSV ve büyük ölçekli ADS-B kayıtlarının okunması,
  \item farklı örnekleme hızındaki telemetri kanallarının zaman ekseninde birleştirilmesi,
  \item hatalı, eksik, donmuş veya fiziksel olarak anlamsız kayıtların temizlenmesi,
  \item uçuş, olay ve zaman satırı seviyesinde etiketlerin yeniden kurulması,
  \item ham sensörlerden fiziksel ve istatistiksel yeni özelliklerin üretilmesi,
  \item klasik ML, derin öğrenme ve fiziksel yöntemlerin uygulanması,
  \item eşik kalibrasyonu, kör test ayrımı ve yanlış alarm analizi,
  \item yanıltıcı başarıların kök neden analizi,
  \item son fizibilite deneyinin tasarlanması ve karar raporunun hazırlanması.
\end{itemize}

Bu nedenle rapor yalnız sonuçları değil, hangi teknik gözlemin bir sonraki
karara yol açtığını da açıklamaktadır.

\section{Kullanılan telemetri ve üretilen genel özellikler}

Veri setlerinde sütun adları ve kayıt biçimleri değişse de çalışılan bilgi
grupları aşağıdaki ortak başlıklarda toplandı.

\begin{longtable}{p{0.22\textwidth}p{0.31\textwidth}p{0.39\textwidth}}
\toprule
\textbf{Bilgi grubu} & \textbf{Ham telemetri örnekleri} & \textbf{Üretilen veya kontrol edilen özellikler} \\
\midrule
\endhead
Zaman ve uçuş kimliği &
Zaman damgası, uçuş/oturum kimliği, mesaj sırası &
Örnekleme aralığı, veri boşluğu süresi, mesaj donması, uçuş bazlı bölme ve olay süresi \\

Konum ve navigasyon &
Enlem, boylam, yerel X--Y--Z, GPS irtifası, uydu sayısı, GPS fix durumu &
İki konum arasındaki mesafe, konumdan hesaplanan hız, GPS hız tutarlılığı,
konum sıçraması, donmuş konum sayacı \\

Hız ve yön &
Yer hızı, kuzey/doğu hızları, dikey hız, heading/yaw &
Hız değişimi, yön değişimi, dönüş oranı, hız bileşenleri arası tutarlılık,
ani hız ve yön sıçramaları \\

Tutum ve dinamik &
Roll, pitch, yaw, açısal hızlar, ivmeler &
Açı farkları, türevler, hareketli ortalama/sapma, dönüş ile yatış açısı
arasındaki fiziksel uyum \\

Otopilot ve kestirim &
İstenen ve ölçülen durumlar, EKF innovation, test ratio, reset ve hata bayrakları &
İstenen--gerçekleşen kontrol hatası, kestirim kalıntıları, sensör güven
göstergeleri ve kalıcı sapma skorları \\

Aktüatör ve enerji &
Motor/PWM komutları, itki komutu, batarya gerilimi/akımı, ESC telemetrisi &
Motor simetrisi, komut--tepki farkı, enerji tüketim değişimi, tek kanalda
donma veya doygunluk \\

Veri kalitesi &
Eksik değer, geçersiz sayı, sentinel değer, kanalın kaybolması &
Eksiklik oranı, kanal kullanılabilirliği, sabitlenme süresi, fiziksel sınır
kontrolü ve değerlendirilebilirlik bayrağı \\
\bottomrule
\end{longtable}

\subsection{Ön işleme ve temizleme yaklaşımı}

Temizleme işlemleri yalnız eksik hücre doldurmaktan ibaret tutulmadı. Aşağıdaki
problemler ayrı ayrı ele alındı:

\begin{itemize}
  \item Aynı fiziksel büyüklüğün farklı veri kaynaklarında farklı sütun
  isimleriyle tutulması eşleştirildi.
  \item Farklı frekansta gelen sensör ve otopilot mesajları ortak zaman çizelgesine
  hizalandı.
  \item GPS olmayan kapalı alan uçuşlarında yerel konum kanalları kullanıldı;
  GPS varmış gibi sahte değer üretilmedi.
  \item Eksik değerler modelin ``eksikliği arıza sanmasına'' yol açmayacak biçimde
  işlendi; ayrıca eksiklik bilgisi ayrı veri kalitesi özelliği olarak korundu.
  \item Donmuş GPS, sıfırda kalan ESC alanları, geçersiz doğruluk değerleri ve
  sentinel değerler tespit edildi.
  \item Bir uçuşta arıza yalnız belirli bir zaman aralığında gerçekleşiyorsa,
  tüm uçuşu arızalı sayan kaba etiketler olay başlangıç ve bitişine göre düzeltildi.
  \item Aynı uçuş veya aynı oturumun eğitim ve test tarafına dağılması engellendi.
\end{itemize}

\subsection{Sentetik anomali kullanımı}

Gerçek arıza örneğinin az olduğu yerlerde hız sapması, konum kayması, donma,
dropout veya yön değişimi gibi kontrollü enjeksiyonlar yalnız değerlendirme
amacıyla kullanıldı. Sentetik veri modelin eğitiminde gerçek arıza varmış gibi
kullanılmadı.

Ayrıca yalnız ``enjeksiyon komutu verildi'' bilgisi yeterli kabul edilmedi.
Değişikliğin ilgili telemetri kanalında gerçekten oluşup oluşmadığı ve modelin
gördüğü özelliklerde gözlenebilir olup olmadığı kontrol edildi. Böylece fiziksel
olarak görünmeyen bir enjeksiyonun kaçırılması model hatası olarak sayılmadı.

\section{Veri kaynakları ve projedeki rolleri}

\begin{longtable}{p{0.20\textwidth}p{0.34\textwidth}p{0.38\textwidth}}
\toprule
\textbf{Veri ailesi} & \textbf{Neden kullanıldı?} & \textbf{Projeye sağladığı temel ders} \\
\midrule
\endhead
ALFA &
Gerçek sabit kanat uçuşları ve aktüatör/arıza örnekleri içerdiği için &
Ham uçuş kayıtlarından daha zengin özellik çıkarıldığında LSTM Autoencoder
gibi yöntemlerin belirli bir platformda güçlü ayrışma sağlayabildiğini gösterdi. \\

UAV Attack &
GPS spoofing, jamming ve hizmet engelleme saldırılarını içerdiği için &
Dosya adı veya senaryo etiketi ile fiziksel telemetri etkisinin aynı şey
olmadığını; bazı saldırıların mevcut sensörlerde görünmeyebildiğini gösterdi. \\

UAV-SEAD &
GPS, irtifa, mekanik ve dış konum anomalileriyle daha çeşitli normal uçuşlar
içerdiği için &
Normal veri miktarını artırmanın her zaman başarı getirmediğini; farklı görev ve
uçuş davranışlarının normal sınıfını zorlaştırdığını gösterdi. \\

RflyMAD &
Gerçek ve simülasyon uçuşlarında arıza zamanları incelenebildiği için &
Tüm uçuşu arızalı kabul eden etiketlerin sonuçları şişirdiğini; gerçek olay
zamanı kullanıldığında operasyonel performansın belirgin şekilde düştüğünü gösterdi. \\

ADS-B &
Çok büyük ölçekte, doğal ve etiketsiz hava trafiği davranışı sağladığı için &
Bir yöntemin sentetik anomalileri yakalamasının yeterli olmadığını; doğal trafik
üzerinde alarm yükünün ayrıca ölçülmesi gerektiğini gösterdi. \\
\bottomrule
\end{longtable}

\section{Proje günlüğü: gözlemden karara ilerleyen süreç}

\subsection{İlk adım: temel anomali modeli}

\textbf{Hedef:} Normal uçuş davranışından uzak satırları otomatik bulmak.

\textbf{Uygulama:} İlk olarak Isolation Forest tabanlı bir yöntem kuruldu.
Bu yöntem etiket gerektirmeden çok boyutlu veride seyrek ve uç değerleri
ayırt edebildiği için uygun bir başlangıç noktasıydı.

\textbf{Karşılaşılan sorun:} Tüm telemetri tek modele verildiğinde yüksek
boyut, eksik değer doldurma ve farklı ölçekler skoru bozdu. Model çoğu zaman
arıza davranışından çok büyük sayısal değere veya veri eksikliğine tepki verdi.

\textbf{Ne yaptım:} Özellikleri navigasyon, dinamik, aktüatör ve veri kalitesi
gibi gruplara ayırdım. Her grup için ayrı skor üretip bunları doğrulama verisine
göre normalize ederek birleştirdim. Değerlendirmeyi yalnız satır seviyesinden
çıkarıp uçuş ve olay seviyesine taşıdım.

\textbf{Sonuç ve düşünce:} Temel ayrışma iyileşti; ancak tek bir genel modelin
farklı uçuş şartlarını güvenilir biçimde kapsamayacağı görüldü. Bu nedenle
zaman bağımlılığını öğrenebilen modellere geçildi.

\subsection{Autoencoder ve LSTM dönemi}

\textbf{Hedef:} Normal telemetrinin çok değişkenli ve zamansal örüntüsünü
öğrenmek; yeniden üretilemeyen bölümleri anomali saymak.

\textbf{Uygulama:} Dense Autoencoder, LSTM Autoencoder ve USAD denendi. LSTM
özellikle uçuş verisinin sıralı yapısını modelleyebildiği için tercih edildi.
Eşik belirlemek için normal doğrulama verisindeki hata dağılımı ve kuyruk
modelleme yöntemleri kullanıldı.

\textbf{Olumlu bulgu:} ALFA verisi ham uçuş kayıtlarıyla büyütülüp fiziksel
özellikler eklendiğinde LSTM Autoencoder yaklaşık 0{,}92 uçuş-seviyesi AUC
değerine ulaştı. Bu, belirli bir platform ve arıza tanımı altında derin
öğrenmenin yararlı sinyal bulabildiğini gösterdi.

\textbf{Karşılaşılan sorun:} Aynı başarı daha heterojen veri setlerine ve farklı
platformlara taşınamadı. Eşik başka platformda doğrudan kullanıldığında dağılım
kayması nedeniyle alarm dengesi bozuldu.

\textbf{Ne yaptım:} Modelin kendisini taşımak ile eşik kalibrasyonunu birbirinden
ayırdım. Yeni platformun yalnız normal uçuşlarıyla eşik yeniden ayarlandı.

\textbf{Sonuç ve düşünce:} Kalibrasyon aktarımı belirgin biçimde iyileştirdi,
ancak platformlar arası genel bir model için yeterli olmadı. Sorunun yalnız
model mimarisi değil, normal davranış dağılımının platforma ve göreve göre
değişmesi olduğu anlaşıldı.

\subsection{Daha fazla veri ve daha fazla özellik}

\textbf{Hedef:} Performans düşüşünün veri azlığından kaynaklanıp kaynaklanmadığını
test etmek.

\textbf{Uygulama:} Normal uçuş sayısı artırıldı; EKF innovation/test-ratio,
motor simetrisi, istenen--ölçülen değer farkları, GPS hız tutarlılığı, donma
sayaçları ve spektral özellikler eklendi.

\textbf{Karşılaşılan sorun:} Sabit kanat verisinde daha fazla ham uçuş yararlı
olurken, çok rotorlu ve görev çeşitliliği yüksek veride daha fazla normal uçuş
modelin öğrenmesi gereken normal davranış uzayını genişletti. Yanlış alarm bazı
eşiklerde azaldı, fakat gerçek arıza yakalama oranı da düştü.

\textbf{Sonuç ve düşünce:} ``Daha fazla veri her zaman daha iyi sonuç verir''
varsayımının bu problem için geçerli olmadığı görüldü. Verinin miktarından önce
görev, platform ve uçuş rejimi bakımından temsil kabiliyeti önemliydi.

\subsection{Süpervizyonlu model ve hazır zaman serisi modeli}

\textbf{Hedef:} Etiketli bölümlerde denetimli öğrenmenin ve büyük ön-eğitimli
zaman serisi modellerinin sınırı aşıp aşamayacağını görmek.

\textbf{Uygulama:} LightGBM ile pencere tabanlı sınıflandırma ve Chronos
zero-shot tahmin modeli denendi. Tahmin edilen sinyal ile ölçülen sinyal
arasındaki fark anomali skoru olarak kullanıldı.

\textbf{Gözlem:} Chronos bazı mekanik arıza örneklerinde önceki yönteme göre
daha iyi yakalama sağladı; ancak düşük yanlış alarm şartında yeterli seviyeye
ulaşamadı. LightGBM ise veri setindeki sınırlı ve dengesiz etiket yapısı nedeniyle
kararlı bir üstünlük göstermedi.

\textbf{Sonuç ve düşünce:} Daha güçlü veya daha güncel bir model kullanmanın,
gözlenebilirlik ve etiket sorununu tek başına çözemeyeceği doğrulandı.

\subsection{Tek bir güçlü özelliğin keşfi}

\textbf{Gözlem:} Özellikleri tek tek incelediğimde, bazı aktüatör arızalarında
itki komutunun tek başına çok güçlü sinyal taşıdığı görüldü. Aynı özellik büyük
bir özellik grubuna eklendiğinde katkısı zayıflıyordu.

\textbf{Ne yaptım:} Bu sinyal için ayrı ve ince bir skorlayıcı oluşturdum.

\textbf{Sonuç:} Kategori içindeki en iyi ayrışmalardan biri elde edildi.
Ancak alarm sistemi seviyesinde diğer kanallarla birleştirildiğinde yanlış alarm
yükü arttı ve genel operasyonel hedef yine karşılanmadı.

\textbf{Ders:} Çok sayıda özellik kullanmak her zaman avantaj değildir. Güçlü
bir sinyal, ilgisiz veya gürültülü sütunlar içinde seyrelerek kaybolabilir.

\subsection{Etiket ve değerlendirme hatalarının düzeltilmesi}

\textbf{Gözlem:} Bazı veri setlerinde ``bu uçuşta bir yerde arıza var'' bilgisi,
tüm uçuş satırlarına arıza etiketi olarak yayılmıştı. Böyle bir değerlendirme,
modelin arıza öncesi normal bölümü de doğru alarm saymasına neden olabiliyordu.

\textbf{Ne yaptım:} Gerçek arıza başlangıç ve bitiş zamanlarını kontrol
kanallarından yeniden çıkardım. Çelişkili arıza kimliklerini karantinaya aldım.
Eğitim, doğrulama ve test ayrımlarını uçuş/oturum kimliği üzerinden kurdum.

\textbf{Sonuç:} Bazı yüksek görünen sonuçlar düştü. Örneğin gerçek olay aralığı
kullanıldığında yaklaşık 0{,}53 yakalama ve saatte 22 civarında yanlış alarm
seviyesi görüldü.

\textbf{Düşünce:} Sonuç düşmüş olsa da bu değişiklik gereklidir. Önceki sayı
ürün performansını değil, kaba etiketle uyumu gösteriyordu. Bundan sonra ana
karar ölçütü olarak olay yakalama ve uçuş saati başına yanlış alarm birlikte
kullanıldı.

\subsection{Derin öğrenmede genlik artefaktının bulunması}

\textbf{Gözlem:} LSTM Autoencoder, Dense Autoencoder ve USAD gibi üç farklı
mimarinin neredeyse aynı uçuşları aynı sırada anomalik görmesi şüpheli bulundu.

\textbf{Ne yaptım:} Eğitilmiş ağın skorunu rastgele ağırlıklı, yani hiç
öğrenmemiş bir ağın skoruyla ve ham sinyal büyüklüğüyle karşılaştırdım.

\textbf{Bulgu:} Sıralama korelasyonu yaklaşık 0{,}96 çıktı. Başka bir ifadeyle,
modelin başarısı büyük ölçüde öğrenilmiş davranıştan değil, büyük genlikli
sinyallerin daha yüksek yeniden üretim hatası oluşturmasından geliyordu.

\textbf{Müdahale:} Kanal başına ölçek etkisi giderildi, robust ölçekleme ve
genlik-normalize skorlar uygulandı.

\textbf{Sonuç:} Görünen başarı büyük ölçüde kayboldu. Bu noktada modelin
öğrendiği sanılan sinyalin önemli bölümünün sayısal artefakt olduğu kanıtlandı.
Sonraki tüm sinir ağı deneylerine ``eğitilmiş model--rastgele model--ham
genlik'' kontrolü zorunlu olarak eklendi.

\subsection{ADS-B ile büyük ölçekli doğal alarm yükü analizi}

\textbf{Hedef:} Etiketli küçük veri setlerinde iyi görünen yöntemlerin, yüz
binlerce uçuş içeren doğal hava trafiğinde ne kadar alarm üreteceğini görmek.

\textbf{Uygulama:} Üç tam günlük ADS-B verisi işlendi. Uçuşlar oluşturuldu;
yer hızı, dikey hız, yön, doğu/kuzey hız bileşenleri, irtifa kaynakları ve
dönüş--yatış tutarlılığı gibi fiziksel kalıntılar üretildi. Veri boşluğu ve
kalite problemleri anomali skorundan ayrı tutuldu.

\textbf{Gözlem:} İlk yaklaşım sentetik anomalilerin büyük bölümünü yakalıyor
görünse de doğal trafikte saatte yaklaşık 25 alarm üretiyordu. Bu yük operasyonel
olarak kabul edilebilir değildi.

\textbf{Ne yaptım:} Bağlama duyarlı LSTM tahmin modeli, fiziksel robust skorlayıcı,
Page CUSUM ve kalibre eşikler karşılaştırıldı. Enjeksiyonun gerçekten gözlenebilir
olup olmadığı ayrıca kontrol edildi.

\textbf{Sonuç:} Fiziksel CUSUM belirli konum sapmalarında anlamlıydı; fakat
farklı anomali profillerine genellenemedi. LSTM genlik testini geçmesine rağmen
çok düşük alarm bütçesinde yeterli olay yakalama oranına ulaşamadı.

\textbf{Düşünce:} Büyük veri, küçük veri üzerindeki başarıyı doğrulamak yerine
doğal davranış çeşitliliğini ve yanlış alarm problemini görünür hâle getirdi.

\section{Kullanılan yöntemler: neden seçildi, ne gösterdi?}

\begin{longtable}{p{0.20\textwidth}p{0.27\textwidth}p{0.34\textwidth}p{0.12\textwidth}}
\toprule
\textbf{Yöntem} & \textbf{Seçilme nedeni} & \textbf{Temel gözlem} & \textbf{Karar} \\
\midrule
\endhead
Isolation Forest &
Etiketsiz veride hızlı ve açıklanabilir başlangıç modeli olması &
Ölçek, eksiklik ve yüksek boyuttan kolay etkilendi; modüler kullanıldığında
iyileşti fakat genelleme sınırlı kaldı &
Referans yöntem \\

Dense Autoencoder &
Zamansal olmayan basit yeniden üretim tabanı sağlamak &
Bazı veri setlerinde LSTM ile aynı sıralamayı üretti; genlik etkisi bulundu &
Tek başına uygun değil \\

LSTM Autoencoder &
Telemetrinin zaman bağımlılığını öğrenmek &
Belirli sabit kanat verisinde güçlüydü; heterojen platformlarda ve artefakt
kontrolünden sonra üstünlük kalıcı olmadı &
Koşullu yararlı \\

USAD &
Adversarial yeniden üretimle daha keskin ayrışma ihtimali &
Ek karmaşıklığa rağmen temel Autoencoder'lardan kararlı biçimde iyi olmadı &
Elenmiştir \\

LightGBM &
Etiketli pencerelerde güçlü tabular sınıflandırma sağlamak &
Etiket azlığı, dengesizlik ve platform farkı nedeniyle yeterli genelleme sağlamadı &
Elenmiştir \\

Chronos &
Ön-eğitimli zaman serisi bilgisini zero-shot kullanmak &
Bazı mekanik olaylarda iyileşme sağladı; düşük alarm yükünde yeterli olmadı &
Araştırma sonucu \\

CUSUM / Page CUSUM &
Küçük fakat kalıcı ortalama değişimlerini biriktirmek &
Konum sapması gibi belirli olaylarda sessiz ve etkiliydi; hızlı veya farklı
arıza tiplerinde kaçırma oranı yükseldi &
Dar kapsamda uygun \\

Fiziksel tutarlılık skorları &
Uçuş kinematiği ve sensörler arası ilişkileri doğrudan kontrol etmek &
Modelden bağımsız, açıklanabilir sinyal sağladı; sensör gözlenebilirliği yoksa
doğal olarak çalışmadı &
Temel güvenlik katmanı \\

Bağlamsal LSTM tahmini &
Uçuş fazı ve geçmiş pencereye göre sonraki değeri tahmin etmek &
Genlik artefaktı kontrolünü geçti; geliştirme ve prova verisi arasında alarm
dengesi kararlı değildi &
Ürün için yetersiz \\
\bottomrule
\end{longtable}

\section{Metrikleri nasıl yorumladım?}

Raporun önceki sürümünde çok sayıda AUC, recall ve yanlış alarm değeri bulunuyordu.
Bu değerler tek başına amaç değil, bir sonraki teknik kararı vermek için kullanıldı.

\subsection{AUC neyi söyledi, neyi söylemedi?}

AUC, anomalik örneklerin normallerden genel olarak daha yüksek skor alıp
almadığını gösterir. Model karşılaştırmak ve bir sinyalin varlığını görmek için
yararlıdır. Ancak tek başına:

\begin{itemize}
  \item hangi eşikte alarm verileceğini,
  \item bir arızanın zamanında yakalanıp yakalanmadığını,
  \item operatöre saatte kaç alarm gideceğini
\end{itemize}

göstermez. Bu nedenle yüksek AUC gördüğümde ``ürün hazır'' sonucu çıkarmadım;
eşik altında olay yakalama ve yanlış alarm yükünü ayrıca kontrol ettim.

\subsection{Recall düştüğünde ne düşündüm?}

Recall, gerçek arızaların ne kadarının yakalandığını gösterir. Recall düşükse
önce şu olasılıklar incelendi:

\begin{enumerate}
  \item Arıza ilgili sensör kanalında gerçekten gözlenebilir mi?
  \item Etiket arızanın gerçek başlangıç ve bitişini doğru gösteriyor mu?
  \item Eşik farklı platform veya görev nedeniyle yanlış mı kalibre edildi?
  \item Normal davranış ile arıza davranışı gerçekten üst üste mi biniyor?
\end{enumerate}

Bu kontrollerden sonra özellik, model veya eşik değişikliğine gidildi.

\subsection{Yanlış alarm yükseldiğinde ne düşündüm?}

Yanlış alarm/saat değeri, sistemin normal uçuşta operatöre ne kadar gereksiz
uyarı üreteceğini gösterir. Eşik gevşetildiğinde recall yükseliyor fakat yanlış
alarm da hızla artıyorsa, model sıralama yapabiliyor olsa bile güvenilir karar
sınırı oluşturamıyor demektir.

Projede tekrar eden ana örüntü tam olarak budur: arızaları yakalayan eşikler
fazla alarm üretmiş; sakin çalışan eşikler ise arızaların önemli bölümünü
kaçırmıştır.

\subsection{Benzer modeller aynı sonucu verdiğinde ne yaptım?}

Birbirinden farklı üç derin öğrenme modelinin aynı örnekleri aynı sırada
anomali görmesi ilk bakışta güven verici gibi durabilir. Ben bunu ayrıca bir
şüphe işareti olarak değerlendirdim. Rastgele model ve ham genlik karşılaştırması
sayesinde ortak sonucun gerçek öğrenmeden değil ölçek etkisinden kaynaklandığı
ortaya çıkarıldı.

\section{Başlıca teknik sorunlar ve uygulanan karşılıklar}

\begin{longtable}{p{0.25\textwidth}p{0.32\textwidth}p{0.35\textwidth}}
\toprule
\textbf{Sorun} & \textbf{Neden önemliydi?} & \textbf{Uyguladığım karşılık} \\
\midrule
\endhead
Veri sızıntısı riski &
Aynı uçuşun eğitim ve testte bulunması sonucu yapay olarak yükseltir &
Uçuş ve oturum bazlı ayrım; normalizasyonun yalnız eğitim verisinde öğrenilmesi \\

Tüm uçuşa yayılan kaba etiket &
Arıza öncesi normal bölüm bile doğru alarm sayılabilir &
Gerçek olay başlangıç/bitişinin kontrol kanallarından çıkarılması \\

Eksik veya donmuş sensör &
Model arıza yerine veri kalitesini öğrenebilir &
Kalite bayrakları, eksiklik oranı, freeze sayaçları ve ayrı kalite katmanı \\

Farklı platform dağılımı &
Aynı hız, titreşim veya komut ölçeği her platformda normal olmayabilir &
Platforma özel normal kalibrasyon ve transfer deneyi \\

Normal uçuş çeşitliliği &
Görev, rüzgâr ve manevra farkı arızaya benzeyebilir &
Uçuş fazı/bağlam özellikleri ve doğal alarm yükü testi \\

Az ve dengesiz arıza örneği &
Denetimli model birkaç senaryoyu ezberleyebilir &
Etiketsiz yöntemler, olay bazlı değerlendirme ve sonuçların güven aralığıyla incelenmesi \\

Genlik baskınlığı &
Autoencoder büyük sayıyı anomali sanabilir &
Robust ölçekleme, kanal normalize hata ve rastgele ağ kontrolü \\

Sentetik anomalinin kolaylığı &
Yapay değişim gerçek arızadan daha görünür olabilir &
Fiziksel anlam, gözlenebilir değişim ve doğal trafik alarm yükü kontrolü \\

Sensör gözlenebilirliği &
Arızanın etkisi kaydedilmiyorsa hiçbir model güvenilir tespit yapamaz &
Önce kanal/instrumentation denetimi; görünmeyen arızanın kapsam dışına alınması \\
\bottomrule
\end{longtable}

\section{Son kontrollü çalışma: GNSS bütünlük tespiti}

Genel amaçlı anomali tespitinin sınırları görüldükten sonra son bir kontrollü
deney tasarlandı. Amaç, problemi daraltarak gerçek uçuşlarda GNSS bütünlük
bozulmasını yakalamanın mümkün olup olmadığını görmekti.

\subsection{Neden motor arızası yerine GNSS seçildi?}

Önce veri enstrümantasyonu incelendi. Gerçek uçuş kayıtlarında ESC alanları
bulunsa da RPM, akım ve sıcaklık değerlerinin tamamının sıfır olduğu görüldü.
Bu nedenle motor arızasını bu kanallarla doğrulamak teknik olarak mümkün değildi.

GNSS tarafında ise konum, hız, doğruluk ve otopilot kestirim bilgileri mevcuttu.
Bu nedenle son deney GNSS bütünlük problemine odaklandı. GPS klasöründeki
kayıtların tamamı otomatik olarak GPS arızası sayılmadı; farklı arıza kimliği
taşıyan altı uçuş karantinaya alındı.

\subsection{Karşılaştırılan üç yaklaşım}

\begin{enumerate}
  \item \textbf{PX4/otopilot yerel göstergeleri:} EKF ve GNSS kalite
  göstergelerinin doğrudan kullanımı.
  \item \textbf{Çok kanallı Page CUSUM:} Konum, hız ve kestirim kalıntılarındaki
  kalıcı değişimin istatistiksel olarak biriktirilmesi.
  \item \textbf{Bağlamsal LSTM:} Normal uçuş geçmişinden bir sonraki telemetri
  durumunun tahmin edilmesi ve tahmin hatasının skorlanması.
\end{enumerate}

\subsection{Ne görüldü?}

\begin{itemize}
  \item Otopilot göstergeleri arızaların önemli bölümüne tepki verdi; fakat
  normal uçuşlarda da saatte onlarca alarm üretti.
  \item CUSUM normal uçuşta daha sakindi. Buna karşılık sıkı eşikte olayları
  kaçırdı; gevşek eşikte bazı veri bölümlerinde başarılı olsa da rüzgâr ve
  simülasyon stresinde uyarı yükü arttı.
  \item LSTM bir prova grubunda yüksek yakalama gösterebildi; ancak geliştirme
  grubunda aynı alarm dengesi korunmadı. Sonuç veri bölümleri arasında kararlı
  değildi.
  \item LSTM'nin genlik artefaktına dayanmadığı ayrıca doğrulandı. Dolayısıyla
  bu son denemedeki problem ``ağ hiçbir şey öğrenmedi'' değil; öğrenilen
  ilişkinin yeni uçuş koşullarına yeterince genellenememesi ve kabul edilebilir
  alarm eşiğinin oluşmamasıdır.
\end{itemize}

\subsection{Neden son kör test açılmadı?}

Geliştirme ve prova verilerinde önceden tanımlanan yakalama--yanlış alarm dengesi
birlikte sağlanamadı. Bu durumda son kör test verisini açmak, modele tekrar tekrar
bakıp test setine uyum sağlama riski doğuracaktı. Bu nedenle son test bölümü
kapalı tutuldu. Bu karar başarısızlığı gizlemek için değil, sonucun metodolojik
bütünlüğünü korumak için alındı.

\section{Nihai fizibilite hükmü}

\subsection{Mevcut kapsam için karar}

\resultno{Mevcut veri setleri, etiket kalitesi ve sensör kapsamıyla genel amaçlı,
platformlar arası çalışan ve operasyonel alarm yükünü karşılayan ML tabanlı İHA
anomali tespit ürünü gerçekleştirilebilir olarak doğrulanmamıştır.}

Bu kararın temel gerekçeleri:

\begin{itemize}
  \item gerçek ve doğru zaman etiketli arıza örneğinin sınırlı olması,
  \item bazı arızaların kaydedilen sensörlerde fiziksel olarak görünmemesi,
  \item platform, görev, rüzgâr ve uçuş fazı değişiminin normal davranışı
  ciddi ölçüde değiştirmesi,
  \item başarılı görünen bazı sonuçların etiket veya ölçek artefaktına dayanması,
  \item gerçek arıza yakalama ile yanlış alarm yükünün aynı anda kabul edilebilir
  seviyeye getirilememesi.
\end{itemize}

\subsection{Neler yapılabilir durumda?}

Çalışma tamamen sonuçsuz değildir. Aşağıdaki daha dar kapsamlar teknik olarak
daha gerçekçidir:

\begin{itemize}
  \item Belirli platform ve belirli arıza ailesi için platforma özel model,
  \item GNSS konum sapması için fiziksel kalıntı + CUSUM tabanlı yardımcı uyarı,
  \item otopilotun mevcut kalite göstergelerini birleştiren açıklanabilir izleme katmanı,
  \item veri kalitesi, donma, dropout ve sensör tutarsızlığı için deterministik kurallar,
  \item yalnız doğru sensörlerle enstrümante edilmiş yeni veri toplama kampanyası
  sonrasında motor/aktüatör özelinde yeni fizibilite çalışması.
\end{itemize}

\subsection{Devam kararı verilirse önerilen yol}

Yeni bir genel model denemek yerine aşağıdaki sıra önerilmektedir:

\begin{enumerate}
  \item Tek bir platform ve tek bir arıza ailesi seçilmelidir.
  \item Arızanın hangi sensörde, hangi gecikmeyle ve ne büyüklükte görüleceği
  fiziksel olarak tanımlanmalıdır.
  \item Arıza başlangıç/bitiş zamanı kontrollü deneyle kaydedilmelidir.
  \item Normal veri; görev, rüzgâr, yük, uçuş fazı ve pilotaj çeşitliliğini
  temsil edecek şekilde toplanmalıdır.
  \item Önce fiziksel kural ve CUSUM referansı kurulmalı, ML yalnız ölçülebilir
  ek katkı sağlıyorsa devreye alınmalıdır.
  \item Başarı AUC ile değil, olay yakalama, tespit gecikmesi ve uçuş saati
  başına yanlış alarm ile kabul edilmelidir.
\end{enumerate}

\section{Yönetim görüşmesinde beklenen sorular ve kısa cevaplar}

\subsection*{``Neden bu kadar farklı model denedin?''}

Her model farklı bir hipotezi test etti. Isolation Forest etiketsiz uç değerleri,
Autoencoder'lar normal örüntünün yeniden üretimini, LSTM zaman bağımlılığını,
LightGBM etiketli ayrımı, Chronos ön-eğitimli tahmin bilgisini, CUSUM kalıcı
küçük sapmaları ve fiziksel skorlar sensörler arası tutarlılığı sınadı. Amaç
rastgele model çoğaltmak değil, sorunun modelden mi, veriden mi, etiketten mi
yoksa gözlenebilirlikten mi kaynaklandığını ayırmaktı.

\subsection*{``LSTM'de iyi değer gördün; neden devam etmedin?''}

Belirli bir veri setinde iyi AUC görüldü. Ancak aynı başarı başka platformda
korunmadı ve bazı Autoencoder sonuçlarının ham sinyal genliğiyle oluştuğu
kanıtlandı. Son bağlamsal LSTM genlik testini geçse bile geliştirme ve prova
uçuşlarında aynı yanlış alarm dengesini sağlayamadı. Bu nedenle tek bir iyi
değer ürün fizibilitesi için yeterli sayılmadı.

\subsection*{``Recall'ı artırmak için eşiği düşürseydin olmaz mıydı?''}

Eşik düşürüldüğünde daha fazla arıza yakalandı, fakat normal uçuşlarda saatte
onlarca alarm oluştu. Böyle bir sistem operatör tarafından kısa sürede
önemsizleştirilir. Güvenlik uygulamasında yakalama oranı ve alarm yükü birlikte
sağlanmalıdır.

\subsection*{``Daha fazla veri neden sorunu çözmedi?''}

Daha fazla veri aynı dağılımdan ve doğru temsil ile gelirse yararlıdır. Burada
ek normal uçuşlar farklı görev, manevra ve koşullar içerdiği için normal davranış
alanını genişletti. Model daha az yanlış alarm üretmeyi öğrenirken bazı gerçek
arıza örneklerini de normal kabul etmeye başladı.

\subsection*{``Sentetik arızalar başarılıysa neden ürün yapılmadı?''}

Sentetik arızalar çoğu zaman gerçek arızadan daha temiz ve görünürdür. Ayrıca
enjeksiyon uygulanmış olması, değişikliğin model girdisinde mutlaka oluştuğu
anlamına gelmez. Bu nedenle sentetik yakalama yalnız yöntem testi olarak
kullanıldı; doğal uçuş alarm yükü ve gerçek olaylar ayrıca değerlendirildi.

\subsection*{``Sorun modelde mi, veride mi?''}

Tek bir neden yoktur. Bazı erken sonuçlarda ölçek ve etiket problemi vardı;
bunlar düzeltildi. Son kontrollü deneyde model gerçekten zamansal ilişki
öğrendi, fakat uçuşlar arası genelleme ve alarm eşiği kararlı olmadı. Daha temel
sınır, arıza örneğinin azlığı, normal davranışın genişliği ve bazı arızaların
mevcut sensörlerde yeterince gözlenememesidir.

\subsection*{``Bu çalışma başarısız mı?''}

Ürün hedefi açısından mevcut kapsam için sonuç \textbf{NO-GO}'dur. Mühendislik
açısından ise neden çalışmadığı ölçülmüş, yanıltıcı başarılar bulunup giderilmiş,
veri ve enstrümantasyon eksikleri belirlenmiş ve hangi dar kapsamların anlamlı
olduğu gösterilmiştir. Bu, yeni yatırım kararının yanlış metriklere dayanmasını
engelleyen somut bir çıktıdır.

\section{Sonuç}

Proje boyunca veri toplama ve işleme hattından model geliştirmeye, fiziksel
özellik üretiminden etiket denetimine, derin öğrenme artefakt analizinden büyük
ölçekli doğal alarm testine kadar bütün zincir ele alındı. İlk sonuçlarla
yetinilmeyip her yüksek değerin kaynağı sorgulandı; sonuçlar düştüğünde bunu
gizlemek yerine neden düştüğü araştırıldı.

Elde edilen nihai teknik sonuç şudur: Mevcut veri ve sensör kapsamı, genel
amaçlı bir İHA anomali tespit ürününü güvenilir biçimde doğrulamak için yeterli
değildir. En doğru devam yolu yeni bir genel model denemek değil; tek platform,
tek arıza ailesi, doğru enstrümantasyon ve kontrollü gerçek arıza zamanlarıyla
dar kapsamlı yeni bir veri kampanyası yürütmektir.

\appendix

\section{Terimlerin kısa karşılığı}

\begin{description}[leftmargin=3.2cm,style=nextline]
  \item[AUC] Modelin anomalik örnekleri normal örneklerden genel olarak ayırma
  gücü. Tek başına operasyonel başarı değildir.
  \item[Recall] Gerçek arıza olaylarının yakalanan oranı.
  \item[Yanlış alarm/saat] Normal uçuş sırasında bir saatte üretilen gereksiz
  alarm sayısı.
  \item[Autoencoder] Normal veriyi sıkıştırıp yeniden üretmeyi öğrenen; yüksek
  yeniden üretim hatasını anomali kabul eden sinir ağı.
  \item[LSTM] Zaman sıralı veride geçmiş bilgiyi kullanabilen sinir ağı yapısı.
  \item[CUSUM] Küçük fakat kalıcı değişimleri zaman içinde biriktiren
  istatistiksel değişim tespit yöntemi.
  \item[Fiziksel kalıntı] İki sensör veya fiziksel ilişki arasındaki beklenen ve
  ölçülen değer farkı.
  \item[Kör test] Model ve eşik kararları tamamlanana kadar açılmayan son
  değerlendirme verisi.
\end{description}

\section{Teknik kanıt dokümanları}

Bu rapor yönetici ve teknik görüşme anlatımı için sadeleştirilmiştir. Ayrıntılı
deney tabloları, karar kayıtları ve son GNSS deneyinin sayısal çıktıları proje
deposunda aşağıdaki belgelerde korunmaktadır:

\begin{itemize}
  \item Ayrıntılı fizibilite raporu:
  \code{docs/final_rapor_ml_fizibilite_2026-07-16.tex}
  \item Geniş teknik rapor:
  \code{docs/ml_anomali_tespiti_fizibilite_final_raporu_2026-07-16.tex}
  \item GNSS son deney raporu:
  \code{artifacts/uav_gnss_integrity_v1/uav_gnss_integrity_v1_final_no_go_report.tex}
  \item GNSS deney yapılandırması:
  \code{configs/uav_gnss_integrity_v1.json}
\end{itemize}

\end{document}
